\section{Realización del Proyecto}
    Para facilitar la carga y posterior aplicación de transformaciones a las imágenes, 
    se decidió usar OpenCV \parencite{opencv_library} para la lectura de las imágenes 
    como arrays de NumPy \parencite{numpy_library}, considerando que OpenCV no lee los 
    canales de un imagen como RGB sino BGR, por lo que se tuvo que realizar esta 
    transformación menor para una adecuada lectura y visualización de las imágenes 
    empleando Matplotlib \parencite{matplotlib_library}.\\

    Para cada una de las imágenes del dataset, se realizó primero una visualización 
    inicial en la que se muestra su versión original, en escala de grises y cada 
    canal de color de forma individual. Esta función permitió dar una inspección 
    sobre qué problema o en qué situación estaba una cierta imagen, e igualmente 
    para determinar si después de aplicar las transformaciones convenientes se 
    alcanzaba una mejora o no visualmente. Adicionalmente, se incluyen los histogramas 
    de las intensidades de los colores por canal con el fin de estimar los parámetros 
    más adecuados para las transformaciones para así lograr una mejora más fina 
    sobre la estética de las imágenes.\\

    Tomando en cuenta esta estimación heurística, se realizó la comparación del antes 
    (imagen original) contra el después de aplicar la transformación, donde de forma 
    general no todas las imágenes mejoraron con una misma técnica, es decir, algunas 
    se les tuvo que aplicar la corrección gamma con un valor negativo mientras que 
    en otras les fue más conveniente usar la transformación gain-bias para lograr un 
    resultado más natural o una mejora en la calidad visual.